% AUTO-GENERATED by slm_harness/scripts/fill_phase2_paper.py. Do not edit.
\newcommand{\floorSummarySentence}{Of the eight subroutines, 3 clear the bar at or below 500M parameters, 1 of them already at 135M; for 3 the rules baseline alone suffices; 2 remain unsolved at every tested size.}
\newcommand{\minSchemaValidityFTPct}{42}
\newcommand{\schemaValiditySentence}{Fine-tuned schema validity is at least 92 percent in 31 of 32 cells; the lone exception is \texttt{path\_normalizer} at 135M with 42 percent. Base instruct models at the same two smallest sizes collapse to zero schema validity on most subroutines under identical prompts, so format reliability is bought by fine-tuning, not by scale.}
\newcommand{\nTrainPerSub}{2000}
\newcommand{\nValPerSub}{150}
\newcommand{\nTestPerSub}{300}
\newcommand{\nTestEval}{250}
\newcommand{\floorTableBody}{\begin{tabular}{lrrrrrll}
\toprule
Subroutine & Rules & 135M & 362M & 494M & 1544M & Floor & Verdict \\
\midrule
\texttt{action\_router} & 0.277 & 0.612 & 0.740 & \textbf{0.936} & \textbf{0.996} & 494M & works at 494M \\
\texttt{evidence\_judge} & 0.983 & 0.540 & 0.484 & 0.844 & 0.932 & -- & rules suffice \\
\texttt{json\_repair} & 0.147 & \textbf{0.924} & \textbf{1.000} & \textbf{0.996} & \textbf{1.000} & 135M & works at 135M \\
\texttt{path\_normalizer} & 0.790 & 0.500 & 0.704 & \textbf{0.956} & \textbf{0.980} & 494M & works at 494M \\
\texttt{read\_span\_selector} & 0.817 & 0.000 & 0.024 & 0.224 & 0.816 & -- & unsolved (best 0.82 at qwen2.5-1.5b) \\
\texttt{search\_hit\_ranker} & 0.993 & 0.160 & 0.176 & 0.948 & 0.996 & -- & rules suffice \\
\texttt{search\_query\_gen} & 0.237 & 0.636 & 0.632 & 0.864 & 0.876 & -- & unsolved (best 0.88 at qwen2.5-1.5b) \\
\texttt{trace\_localizer} & 1.000 & 0.812 & 0.936 & 1.000 & 1.000 & -- & rules suffice \\
\bottomrule
\end{tabular}}
\newcommand{\floorTableNarrative}{Of the eight subroutines, 3 clear the bar at or below 500M parameters, 1 of them already at 135M; for 3 the rules baseline alone suffices; 2 remain unsolved at every tested size. Bold cells clear the 0.90 bar while beating rules by two points. At every declared floor the fine-tuned specialist beats the rules baseline on a paired exact McNemar test (p < 0.001, identical test items).}
\newcommand{\ftVsBaseNarrative}{At 135M parameters the fine-tuned specialists average 0.523 success across subroutines while the identical base instruct model averages 0.030 under the same prompts and retry budget, with mean schema validity of only 0.031. At 1544M the gap is 0.950 against 0.448. Fine-tuned schema validity is at least 92 percent in 31 of 32 cells; the lone exception is \texttt{path\_normalizer} at 135M with 42 percent. Base instruct models at the same two smallest sizes collapse to zero schema validity on most subroutines under identical prompts, so format reliability is bought by fine-tuning, not by scale. This replicates the co-design lesson of the previous study at one tenth to one hundredth the previous deployment size.}
\newcommand{\rulesSufficeNarrative}{For \texttt{trace\_localizer}, \texttt{search\_hit\_ranker} and \texttt{evidence\_judge}, a deterministic solver we wrote in an afternoon already clears the reliability bar (1.000, 0.993, 0.983 respectively), so the honest engineering verdict there is to ship the rules, not a model. At the other end, rules score only 0.277, 0.237, 0.147 on \texttt{action\_router}, \texttt{search\_query\_gen} and \texttt{json\_repair}, and these are exactly the subroutines where learned specialists earn their place.}
\newcommand{\costNarrative}{Training all 32 checkpoints took 59 GPU minutes and the full evaluation grid took 7 GPU minutes, about \$3.53 at the \$3.19 hourly rate of the H100 NVL used. At the floor sizes, batched cost per successful call is small enough to be effectively free next to orchestrator tokens. For example \texttt{path\_normalizer} at 494M costs \$0.000003 per success (0.00s per call); \texttt{action\_router} at 494M costs \$0.000004 per success (0.00s per call); \texttt{json\_repair} at 135M costs \$0.000061 per success (0.06s per call). These figures are batched-throughput costs on the rental GPU; a 135M or 360M specialist also runs locally on CPU, where the marginal cost is zero.}
\newcommand{\failureNarrative}{Where the bar is missed, the shortfall is concentrated rather than uniform. \texttt{evidence\_judge} peaks at 0.932 (qwen2.5-1.5b); \texttt{read\_span\_selector} peaks at 0.816 (qwen2.5-1.5b); \texttt{search\_hit\_ranker} peaks at 0.996 (qwen2.5-1.5b); \texttt{search\_query\_gen} peaks at 0.876 (qwen2.5-1.5b); \texttt{trace\_localizer} peaks at 1.000 (qwen2.5-0.5b). Success is not perfectly monotone in size for \texttt{evidence\_judge}, consistent with the single un-tuned hyperparameter schedule shared across all cells. The largest size-to-size gains fall exactly on the SmolLM2-to-Qwen family boundary for \texttt{evidence\_judge} (0.484 to 0.844), \texttt{search\_hit\_ranker} (0.176 to 0.948), so family and scale are confounded there and the 360M-to-500M step should not be read as a pure scale effect. Residual errors at the floor sizes are dominated by near-miss spans and confusable candidates rather than schema violations, which the guards catch at deployment time.}
\newcommand{\conclusionNarrative}{A deterministic, schema-verifying harness moves the deployable size of real developer-agent subroutines well below one billion parameters. Of the eight subroutines, 3 clear the bar at or below 500M parameters, 1 of them already at 135M; for 3 the rules baseline alone suffices; 2 remain unsolved at every tested size. The recipe is unchanged from our previous study, narrow the task, externalize state, verify deterministically, and fine-tune the smallest model that clears the bar, applied one order of magnitude lower in scale. The parameter-floor map, not any single accuracy number, is the deliverable. It tells a harness designer which sub-tasks to hand to a local 135M to 500M specialist, which to leave to rules, and which still need a larger model behind an MCP boundary.}
